\section{Related Work}
\label{sec:related_work}

\subsection{Vision-Language-Action Models}

VLAs repurpose pretrained vision-language models as robotic policies by fine-tuning them on demonstration data.
RT-2~\cite{brohan2023_rt2} established the paradigm by showing that a large vision-language model can be directly fine-tuned to output robotic actions tokenized as text, transferring web-scale knowledge to manipulation tasks.
OpenVLA~\cite{kim2024_openvla} made this architecture open-source, building on the Prismatic VLM framework~\cite{karamcheti2024} with a frozen SigLIP encoder~\cite{zhai2023_siglip} and a Llama-2 backbone.
$\pi_0$~\cite{black2024_pi0} departed from autoregressive action decoding in favor of flow matching, while retaining the frozen-encoder-plus-LLM structure.
A common architectural choice across these models is to freeze the vision encoder and fine-tune only the LLM and projector, embedding an implicit assumption: fine-tuning injects action-relevant knowledge into the LLM.
Our work tests this assumption directly at the representational level and finds that it does not hold---action-relevant geometric structure originates in the vision encoder and is routed, not created, by fine-tuning.

\subsection{Understanding VLA Representations}

A growing body of work investigates how VLAs process information internally.
Haon et al.~\cite{haon2025} performed the first mechanistic interpretability study of VLAs, identifying causally relevant attention heads and demonstrating activation steering for behavior modification.
Grant et al.~\cite{grant2026} showed across six VLA architectures that the visual pathway dominates action generation, with language stream ablations causing minimal behavioral degradation---a behavioral-level finding consistent with our representational-level observation of vision encoder dominance.
Swann et al.~\cite{swann2026} applied sparse autoencoders to OpenVLA, discovering interpretable monosemantic features at the individual-neuron level.
Rizzardo et al.~\cite{rizzardo2026} characterized VLA representations through the lens of observability and controllability, measuring how well internal features predict and steer downstream behavior.
Molinari et al.~\cite{molinari2025} probed OpenVLA for world-model properties, finding that the model encodes spatial and physical state information in its intermediate representations.

These studies illuminate different facets of VLA internals but leave a gap that our work addresses.
Grant et al.\ operate at the behavioral level (input ablations and output measurements) without examining representation geometry.
Swann et al.\ and Molinari et al.\ analyze individual features or semantic properties without quantifying how much action-relevant structure each architectural component contributes.
Haon et al.\ identify causal mechanisms but do not decompose where the underlying representations originate.
No prior work compares trained and untrained models to isolate training's contribution, and none reports the geometric-functional dissociation we observe between vision-language and vision-only encoders.

\subsection{Representational Analysis Methods}

Our methodology draws on a toolkit developed across neuroscience and deep learning.
Representational Similarity Analysis (RSA)~\cite{kriegeskorte2008}, originally designed to compare neural recordings with computational models, quantifies geometric alignment between representation spaces via pairwise dissimilarity correlations.
Centered Kernel Alignment (CKA)~\cite{kornblith2019} provides a complementary global measure invariant to orthogonal transformations, enabling comparison of representations across architectures and dimensionalities.
Linear probing~\cite{alain2017} assesses the functional content of representations by training lightweight decoders on frozen activations; Belinkov~\cite{belinkov2022} surveys its extensive use in NLP for analyzing linguistic structure at each layer of pretrained models.

These tools have been widely applied to vision models and language models in isolation, but their joint application to VLA pipelines---where vision, language, and action spaces intersect---is largely unexplored.
We combine RSA, CKA, and probing in a factorial design to provide the first systematic representational decomposition of a VLA, measuring both geometric alignment and functional decodability across all architectural components and training conditions.
The multi-metric approach proves essential: we show that linear CKA and RSA can produce divergent rankings of the same representations (Section~\ref{sec:vision_dominance}), underscoring the need for methodological plurality in representation analysis.

\subsection{Vision-Language Pretraining}

The vision encoders in modern VLAs inherit their representations from large-scale pretraining.
CLIP~\cite{radford2021_clip} and SigLIP~\cite{zhai2023_siglip} learn visual features aligned to language through contrastive objectives on image-text pairs, while DINOv2~\cite{oquab2024_dinov2} learns from self-supervised vision objectives without language supervision.
Prior work has shown that these pretraining paradigms produce qualitatively different representation geometries, but their downstream consequences for action-relevant structure in VLAs have not been characterized.
Our finding of a geometric-functional dissociation---VL-pretrained encoders achieve substantially higher geometric alignment with the action space than DINOv2 yet comparable functional decodability---provides the first direct evidence that language-grounded pretraining shapes the geometry, but not the information content, of action-relevant visual representations.
